\subsection*{定理 5.6.8.}
\underline{设 $X$ 为代数闭域 $k$ 上的 $n$ 维光滑概形, 并设}
\[
\rho:X\longrightarrow\mathbb P^1_k
\]
\underline{为一般纤维光滑的射影态射. 假设:}

\smallskip
\noindent
\textup{(5.6.8.1)}\quad
\underline{几何一般纤维配备相对丰沛可逆层 $\underline{\mathcal O}_X(1)$ 的限制后, 满足 (LV);}

\smallskip
\noindent
\textup{(5.6.8.2)}\quad
\underline{当 $i\neq n$ 时映射 (5.3.2) 为同构, 当 $i=n$ 时为满射 (特别地, 当 $\rho$ 来自 Lefschetz 线束时此条件成立, 参见 5.3.4). 则 Leray 谱序列}
\[
E_2^{p,q}=H^p(\mathbb P^1,R^q\rho_*\mathbb Q_\ell)
\Longrightarrow H^{p+q}(X,\mathbb Q_\ell)
\]
\underline{退化.}

\smallskip
\noindent
\underline{注.}
将 $\mathbb P^1_k$ 换成任意光滑、不可约射影曲线, 同一定理仍成立, 证明也相同.

\underline{证明.}
证明中我们将层 $\mathbb Q_\ell(1)$ 与 $\mathbb Q_\ell$ 等同; 由于 $k$ 代数闭, 这是可行的. 显然可以假设 $X$ 连通且 $\rho$ 为\underline{满射}, 因而为平坦态射, 纤维维数为 $n-1$. 设
\[
\xi_X\in H^2(X,\mathbb Q_\ell)
\]
为 $\underline{\mathcal O}_X(1)$ 的第一 Chern 类, 并设
\[
\nu:U\hookrightarrow\mathbb P^1
\]
为一个开子集的包含映射, 使 $\rho$ 在此开子集上光滑. 由 (5.6.8.1) 和引理 5.2.6, 杯积映射
\[
\xi_X^i:
\nu^*R^{n-1-i}\rho_*\mathbb Q_\ell
\longrightarrow
\nu^*R^{n-1+i}\rho_*\mathbb Q_\ell
\]
在 $i\geq1$ 时为同构. 应用 $\nu_*$ 并利用 (5.6.8.2), 当 $i\geq2$ 时得到同构
\begin{equation}\tag{5.6.8.3}
\xi_X^i:R^{n-1-i}\rho_*\mathbb Q_\ell
\overset{\sim}{\longrightarrow}
R^{n-1+i}\rho_*\mathbb Q_\ell,
\end{equation}
而当 $i=1$ 时得到同构
\begin{equation}\tag{5.6.8.4}
\xi_X:R^{n-2}\rho_*\mathbb Q_\ell
\overset{\sim}{\longrightarrow}
\nu_*\nu^*R^n\rho_*\mathbb Q_\ell.
\end{equation}
后者分解为
\begin{equation}\tag{5.6.8.5}
R^{n-2}\rho_*\mathbb Q_\ell
\xrightarrow{\ \xi_X\ }
R^n\rho_*\mathbb Q_\ell
\xrightarrow{\ \text{伴随}\ }
\nu_*\nu^*R^n\rho_*\mathbb Q_\ell.
\end{equation}
因而有典范分解
\begin{equation}\tag{5.6.8.6}
R^n\rho_*\mathbb Q_\ell
\xleftarrow[\sim]{\ \xi_X+\mathrm{incl}\ }
R^{n-2}\rho_*\mathbb Q_\ell\oplus
\operatorname{Local}^n\mathbb Q_\ell,
\end{equation}
其中
\begin{equation}\tag{5.6.8.7}
\operatorname{Local}^n\mathbb Q_\ell
=\operatorname{Ker}\!\left(
R^n\rho_*\mathbb Q_\ell
\longrightarrow\nu_*\nu^*R^n\rho_*\mathbb Q_\ell
\right).
\end{equation}
层 $\operatorname{Local}^n\mathbb Q_\ell$ 集中于 $\mathbb P^1$ 上 $\rho$ 的纤维不光滑的点. 将此事实与 (5.6.8.3) 结合, 得
\begin{equation}\tag{5.6.8.8}
\xi_X^{i+1}:
H^j(\mathbb P^1,R^{n-2-i}\rho_*\mathbb Q_\ell)
\overset{\sim}{\longrightarrow}
H^j(\mathbb P^1,R^{n+i}\rho_*\mathbb Q_\ell),
\qquad i\geq0,\quad j\geq1.
\end{equation}
\footnote{印本的 (5.6.8.8) 中次数移动为奇数, 映射仅标为 $\xi_X$, 且从 $i=1$ 开始. (5.6.8.3)、(5.6.8.6)、下一个交换方块中的下方箭头, 以及稍后特取的 $E_2^{1,n-2}\simeq E_2^{1,n}$, 都要求此处给出的改正形式.}

现在证明退化性. 由于除 $p=0,1,2$ 外都有 $E_2^{p,q}=0$, 只须证明所有微分
\[
d_2:H^0(\mathbb P^1,R^i\rho_*\mathbb Q_\ell)
\longrightarrow H^2(\mathbb P^1,R^{i-1}\rho_*\mathbb Q_\ell)
\]
均为零. 当 $q\leq n-1$ 时, 置
\begin{equation}\tag{5.6.8.9}
\operatorname{Prim}(E_2^{p,q})
=\operatorname{Ker}\!\left(
\xi_X^{n-q}:E_2^{p,q}\longrightarrow E_2^{p,2n-q}
\right).
\end{equation}
公式 (5.6.8.3) 和 (5.6.8.6), 特别是下式的单射性:
\[
\xi_X:H^p(\mathbb P^1,R^{n-2}\rho_*\mathbb Q_\ell)
\longrightarrow H^p(\mathbb P^1,R^n\rho_*\mathbb Q_\ell),
\]
给出 $q\leq n-1$ 时的分解
\begin{equation}\tag{5.6.8.10}
E_2^{p,q}
\xleftarrow[\sim]{\ 1+\xi_X+\xi_X^2+\cdots\ }
\operatorname{Prim}(E_2^{p,q})
\oplus\operatorname{Prim}(E_2^{p,q-2})
\oplus\operatorname{Prim}(E_2^{p,q-4})
\oplus\cdots.
\end{equation}
\footnote{印本只列出前三个本原直和项, 而其箭头标记已含省略号; 末尾的 $\oplus\cdots$ 将同一个有限分解明确写出.}
